% ====================== PAQUETES Y CONFIGURACIONES ===================== %
% Formatos y esquema del documento
\usepackage[letterpaper,left=1.2in,right=1in,top=1in,bottom=1in]{geometry}
\usepackage[utf8]{inputenc}
\usepackage[spanish,es-tabla]{babel}
\usepackage[explicit]{titlesec}  % Títulos y formatos
\usepackage[titles]{tocloft}     % Tabla de contenidos
\usepackage{caption,subcaption}
\usepackage{fancyhdr,listings,adjustbox,setspace,ragged2e}
\usepackage[title]{appendix}
\renewcommand{\baselinestretch}{1.2}

% Fórmulas Matemáticas
\usepackage{amsmath,amssymb,amsfonts,mathrsfs,esint,cancel,latexsym,mathtools,bm,scalerel,nicematrix,amsthm}
\setcounter{MaxMatrixCols}{18}

% Gráficos y Figuras
\usepackage{float,color,xcolor,graphics,graphicx,pgf,tikz,pgfplots,pstricks,pstricks-add,epsfig}
\pgfplotsset{compat=1.16}

% Texto, tablas y marcos de texto
\usepackage{enumerate,lipsum,ulem,enumitem,setspace,multirow,soul,verbatim,fancybox,array}
\usepackage[most]{tcolorbox}
\usepackage{varwidth}
\usepackage[explicit]{titlesec}
\usepackage{microtype}


% Bibliografía y Referencias
%\usepackage{cite}
\definecolor{myblues}{rgb}{0.00, 0.18, 0.39}
\usepackage[colorlinks=true,linkcolor=black,citecolor=myblues]{hyperref}
\usepackage[square,numbers]{natbib}
\bibliographystyle{IEEEtranN}


% Comandos propios y teoremas/lemas
\newcommand{\oequal}[1]{\ensuremath{\overset{(#1)}{=}}}
\newcommand{\mean}[1]{\mathcal{E}\left\{ #1 \right\}}
\newcommand\norm[1]{\left\lVert#1\right\rVert}
\providecommand{\abs}[1]{\lvert#1\rvert}
\numberwithin{equation}{chapter}
%\newtheorem{lemma}{Lema}
%\newtheorem{thm}{Teorema}
%\newtheorem{example}{Ejemplo}
%\newtheorem{remark}{Observación}
%\renewcommand{\chaptertitlename}{Capítulo}
%\renewcommand{\proofname}{Demostración}

